% ===== cie-cuoti 模板：错题重做卷（正面题目 / 反面官方 Mark Scheme）=====
% 用法：复制本文件与 preamble.tex 到工作区，替换 <<...>> 占位后 xelatex 编译两遍。
% ⚠ 填充文本中的特殊字符必须转义：文件名里的下划线写 \_（如 蜜蜂试卷\_XXXXXX.pdf），& % # $ 同理。
% 结构：封面(含映射表) → \qfront{...} → \msback{...} → …… 双面打印后题与 MS 一一对应。
\documentclass[11pt]{article}
% ===== cie-cuoti 公共导言区（被 template_redo.tex / template_notebook.tex \input）=====
\usepackage[UTF8]{ctex}
\usepackage[a4paper,margin=1.6cm,top=1.9cm,bottom=1.7cm,headsep=6mm]{geometry}
\usepackage{amsmath,amssymb}
\usepackage{enumitem}
\usepackage{booktabs}
\usepackage{array}
\usepackage{xcolor}
\usepackage{tcolorbox}
\tcbuselibrary{skins,breakable}
\usepackage{fancyhdr}
\usepackage[hidelinks]{hyperref}

\definecolor{mainc}{RGB}{30,90,160}
\definecolor{accc}{RGB}{200,70,40}
\definecolor{okc}{RGB}{0,120,80}
\definecolor{lightgray}{RGB}{245,245,245}

\newcommand{\dd}{\displaystyle}
\setlist[enumerate]{leftmargin=2.2em,itemsep=1pt,topsep=2pt,parsep=0pt}
\setlist[itemize]{leftmargin=1.6em,itemsep=1pt,topsep=2pt,parsep=0pt}
          % 即 assets/preamble.tex，同目录复制
\usepackage{tabularx}

\definecolor{qcol}{RGB}{25,80,150}
\definecolor{mscol}{RGB}{0,110,70}
\definecolor{warn}{RGB}{190,50,40}

% ---------- 正面：题目（无答案）----------
% #1 序号  #2 练习册题号  #3 代码/官方试卷/分值行  #4 题干  #5 选项
\newcommand{\qfront}[5]{%
\thispagestyle{fancy}%
\begin{tcolorbox}[enhanced,colback=qcol!8,colframe=qcol,boxrule=1pt,arc=1.5mm,
  title={\bfseries\color{white}重做练习 \hspace{2mm}#1\ /\ \TOTALQ\hspace{6mm}练习册 #2},colbacktitle=qcol]
{\small #3}
\end{tcolorbox}\vspace{2mm}
\begin{tcolorbox}[enhanced,colback=white,colframe=qcol!50,boxrule=0.6pt,arc=1.2mm,left=3mm,right=3mm,top=2mm,bottom=2mm]
#4
\end{tcolorbox}\vspace{2mm}
\begin{tcolorbox}[enhanced,colback=white,colframe=qcol!50,boxrule=0.6pt,arc=1.2mm,left=3mm,right=3mm,top=2mm,bottom=2mm]
\textbf{选项}\\[1mm]#5
\end{tcolorbox}\vspace{3mm}
\begin{tcolorbox}[enhanced,height fill,colback=lightgray,colframe=gray!40,boxrule=0.5pt,arc=1.2mm,left=2mm,right=2mm,top=1.5mm,bottom=1.5mm,valign=top]
\textbf{我的作答：}\underline{\hspace{18mm}}\hfill\textbf{所用时间：}\underline{\hspace{18mm}}\\[2mm]
\textbf{计算过程：}
\end{tcolorbox}
\clearpage}

% ---------- 反面：官方 Mark Scheme ----------
% #1 序号  #2 练习册题号  #3 官方出处引用+文件名  #4 正确答案  #5 评分要点
\newcommand{\msback}[5]{%
\thispagestyle{fancy}%
\begin{tcolorbox}[enhanced,colback=mscol!8,colframe=mscol,boxrule=1pt,arc=1.5mm,
  title={\bfseries\color{white}官方 Mark Scheme\hspace{6mm}对应：重做练习 #1（练习册 #2）},colbacktitle=mscol]
{}
\end{tcolorbox}\vspace{2mm}
\begin{tcolorbox}[enhanced,colback=white,colframe=mscol!55,boxrule=0.6pt,arc=1.2mm,left=3mm,right=3mm,top=2mm,bottom=2mm]
\textbf{官方出处与原文}\\[1mm]
{\small #3}\\[2.5mm]
\textbf{正确答案：}{\large\bfseries\color{mscol}#4}
\end{tcolorbox}\vspace{2mm}
\begin{tcolorbox}[enhanced,colback=white,colframe=mscol!55,boxrule=0.6pt,arc=1.2mm,left=3mm,right=3mm,top=2mm,bottom=2mm]
\textbf{评分要点}\\[1mm]#5
\end{tcolorbox}
\clearpage}

% MS 无法获取时，用这个替代 \msback：
\newcommand{\msmissing}[3]{%  #1 序号 #2 练习册题号 #3 已尝试的检索路径说明
\thispagestyle{fancy}%
\begin{tcolorbox}[enhanced,colback=accc!8,colframe=accc,boxrule=1pt,arc=1.5mm,
  title={\bfseries\color{white}Mark Scheme 缺失\hspace{6mm}对应：重做练习 #1（练习册 #2）},colbacktitle=accc]
{}\end{tcolorbox}\vspace{2mm}
\begin{tcolorbox}[enhanced,colback=white,colframe=accc!55,boxrule=0.6pt,arc=1.2mm,left=3mm,right=3mm,top=2mm,bottom=2mm]
\textbf{未找到官方 Mark Scheme}\\[1mm]
已尝试：#3\\[1mm]
本题不提供答案与评分标准，请人工核对或提供官方 MS 后补录。
\end{tcolorbox}\clearpage}

\newcommand{\TOTALQ}{<<总题数>>}   % ← 替换为实际数字，如 16

\begin{document}
\pagestyle{fancy}\fancyhf{}
\fancyhead[L]{\small\itshape <<科目代码>> <<科目名>> · 错题重做卷}
\fancyhead[R]{\small <<来源文件名>>}
\fancyfoot[C]{\small\thepage}
\renewcommand{\headrulewidth}{0.4pt}

% ================= 封面 =================
\begin{center}
\vspace*{18mm}
  {\Huge\bfseries\color{qcol} CIE <<科目代码>> 错题重做练习卷}\\[6mm]
  {\large <<主题范围>>（共 \TOTALQ\ 题）}
\end{center}
\vspace{6mm}
\begin{tcolorbox}[enhanced,colback=white,colframe=qcol!60,boxrule=0.8pt,arc=1.5mm,left=4mm,right=4mm,top=2.5mm,bottom=2.5mm]
\textbf{使用说明}
\begin{itemize}
  \item 本卷仅收录卷面判定为\textbf{做错}的题目；做对的题未混入。
  \item 每页正面为一道题的题干与选项（无答案），反面为对应的\textbf{官方 Mark Scheme}；双面打印时正反面一一对应。
  \item 所有 Mark Scheme 均来自官方文件（本地缓存 \texttt{$\sim$/Documents/CIE-ms/}）。
\end{itemize}
\end{tcolorbox}
\vspace{4mm}
\begin{center}
\renewcommand{\arraystretch}{1.3}
{\small
\begin{tabular}{clll}
\textbf{重做} & \textbf{练习册题号} & \textbf{官方试卷} & \textbf{考季}\\
% ← 每题一行：1 & Q26 & 9701/11 & 2015 May/June \\
<<映射表行>>
\end{tabular}}
\end{center}
\clearpage

% ================= 题目 × N =================
% 示例（务必删除注释后再编译）：
%
% \qfront{1}{Q26}{卷面代码 \texttt{s15-11-3}\quad|\quad 9701/11 Paper 1 Multiple Choice (AS), May/June 2015, Q3\quad|\quad 分值：1 分}{
% $1.00\ \mathrm{g}$ of carbon is combusted in a limited supply of pure oxygen \ldots
% }{A.\ $1\ \mathrm{dm^3}$\qquad B.\ $1.5\ \mathrm{dm^3}$\qquad C.\ $2\ \mathrm{dm^3}$\qquad D.\ $3\ \mathrm{dm^3}$}
% \msback{1}{Q26}{Cambridge International AS \& A Level Chemistry 9701/11, Paper 1 Multiple Choice, May/June 2015, Mark Scheme (Published):\ \textbf{Question 3 \quad Answer A \quad Marks 1}\par\smallskip{\footnotesize 文件：\texttt{9701\_s15\_ms\_11.pdf}（本地缓存，官方）}}{A（$1\ \mathrm{dm^3}$）}{
% 评分要点正文 \ldots\par\smallskip\textcolor{warn}{\textbf{⚠ 红笔订正与官方不一致时的警示写在这里}}}

<<全部 qfront/msback 配对>>

\end{document}
